\documentclass[11pt]{article}
\usepackage[margin=1.05in]{geometry}
\usepackage[T1]{fontenc}
\usepackage{amsmath,amssymb,amsthm,mathtools}
\usepackage{booktabs,longtable,array,graphicx}
\usepackage[expansion=false]{microtype}
\usepackage{xcolor}
\usepackage[colorlinks=true,linkcolor=blue!60!black,citecolor=blue!60!black,urlcolor=blue!60!black]{hyperref}
\graphicspath{{../results/}}

% Status macros (same conventions as the execution-plan document).
\newcommand{\statusV}{\textsf{\footnotesize[V]}}
\newcommand{\statusM}{\textsf{\footnotesize[M]}}
\newcommand{\statusP}{\textsf{\footnotesize[P]}}
\newcommand{\statusS}{\textsf{\footnotesize[S]}}
\newcommand{\statusC}{\textsf{\footnotesize[C]}}
\newcommand{\statusX}{\textsf{\footnotesize[X]}}
\newcommand{\ORS}{\mathsf{ORS}}
\newcommand{\Otilde}{\widetilde{O}}
\newtheorem{problem}{Problem}
\newtheorem{remark}{Remark}

\title{A Randomness Audit of the Polylog-/ORS-Update\\ Dynamic Matching Pipelines\\[2pt]
{\large Task A0 technical note, v0.3 --- companion to the execution plan and to \texttt{audit\_lab}}}
\author{Prepared for Ruifeng Cao}
\date{June 11, 2026}

\begin{document}
\maketitle

\begin{abstract}
\noindent Task A0 of the execution plan asks: for every random choice inside the
pipelines BKSW23, Beh23, ABR24, BG24, AKK25, determine (i) what the randomness
buys, (ii) the failure mode if it is fixed, and (iii) whether limited
independence suffices. This note records read-through passes~1--2 (arXiv full texts
where retrievable, June 2026), the resulting machine-readable audit table
(Table~\ref{tab:audit}, auto-generated from \texttt{auditlib/audit\_table.py}),
and the experimental evidence E1--E5. The headline structural finding: once
the already-deterministic layers are inventoried (boosting is
\emph{derandomized} by Tirodkar; greedy, ORS-charging, and bookkeeping carry
no coins), the entire derandomization problem for the ORS-parameterized
pipelines reduces to exactly \emph{three} modules --- the opportunistic
sampling matcher, the Behnezhad-21 estimation interface, and randomized vertex
sparsification. The estimation interface provably admits no deterministic
drop-in (plan Lemma~3.6 / Cor.~3.7); the sampling matcher is a \emph{search
problem with a promise}, outside that lower bound, and is therefore the
genuinely open algorithmic target (refined Task A5). Status tags
\statusV/\statusM/\statusP/\statusS\ follow the plan.
\end{abstract}

\section{Scope, sources, and verification status}\label{sec:scope}

The five pipelines under audit and the sources actually read in this pass:

\begin{center}\footnotesize
\begin{tabular}{@{}llll@{}}
\toprule
Paper & arXiv source & Read in this pass & Residue (\statusM) \\
\midrule
BG24~\cite{BG24}   & 2404.06069v4 & \S1--\S4.4 (main algorithm, full) & \S5 (ORS upper bd.; not audit-relevant) \\
AKK25~\cite{AKK25} & 2406.13573   & \S1--\S4.1 + Lem.~2.5 proof       & \S4.2 recursion, \S5 (overview only) \\
BKSW23~\cite{BKSW23}& 2207.07438v3 & abstract; App.~A Prop.~A.1; JACM \S2/\S5 & \S5+ robustness lemmas \\
Beh23~\cite{Beh23}  & 2207.07607v2 & abstract + \S1 technique overview & body (EDCS characterization etc.) \\
ABR24~\cite{ABR24}  & 2307.08772v1 & \S1--\S4 (Thm~1: adaptive, worst-case \statusV) & \S5 dynamic details \\
Beh21~\cite{Beh21}  & 2106.02942v3 & \S1--\S3 + \S5 statements & \S3.1--3.2 proof internals \\
\bottomrule
\end{tabular}
\end{center}

\noindent Also verified at statement level: Kiss'22 vertex sparsification in
its operative form, BKSW23 Prop.~A.1 (``a randomized algorithm which for each
update to $G$ makes an update to $O(\log^2 n/\varepsilon^3)$ contracted
subgraphs, such that w.h.p.\ \dots''); Tirodkar's derandomization of the
McGregor boosting framework as stated in AKK25 (Prop.~2.2); and
\cite{ACK19} Thm.~6 ($\Omega(n^2)$ queries to \emph{find} a maximal
matching, even randomized). Remaining \statusM\ residue: BDHSS19/Sol16
maintenance internals, BKSW23 \S5+ robustness lemmas, Kis22 full internals.

\begin{remark}[Version numbering]
AKK25 cites ``[BG24, Lemma 9]'' and ``[BG24, Lemma 14]''. In the v4 text
those statements are \emph{Lemma 7} (RandomSampling) and \emph{Lemma 11}
($\sum_i 1/d_i \le O(\alpha^{-2}\,\ORS_n(r(1{-}\alpha))\log n)$). All BG24
references below are to v4.
\end{remark}

\section{Pipeline anatomy: where the coins actually live}\label{sec:anatomy}

Read-through pass 1 supports the following layered decomposition. The
right-hand column is the verified location of the claim.

\begin{center}\footnotesize
\begin{tabular}{@{}p{0.30\textwidth}cp{0.40\textwidth}@{}}
\toprule
Layer & Coins? & Verified location \\
\midrule
Approximation boosting (McGregor framework) & \textbf{no} & AKK25 Prop.~2.2: ``de-randomized in [Tir18]'' \\
Greedy matching on explicit edge lists & no & AKK25 Fact~2.1; BG24 \textsc{GreedyMatching} (\S4.1) \\
ORS charging / amortization & no & BG24 Lem.~11; AKK25 Lem.~2.5 (its probabilistic method is \emph{analysis-only}) \\
Certificate / batch bookkeeping ($H_{\mathrm{cert}}$; $G_{\mathrm{old}},G_{\mathrm{batch}},G_{\mathrm{match}}$) & no & BG24 Alg.~2; AKK25 Alg.~2 + Fig.~1 \\
Bounded-degree $(1+\varepsilon)$ maintenance and the periodic-recomputation paradigm (Gupta--Peng) & no & deterministic; BLM20 Lem.~6.1; BKSW23 JACM \S2 \\
\midrule
Opportunistic sampling matcher & \textbf{yes} & BG24 Alg.~1/Lem.~7 + Claim~8; AKK25 Alg.~1/Lem.~3.1 + Claims~3.3--3.4 \\
Sublinear size-estimation interface (Beh21) & \textbf{yes} & BG24 Prop.~4; AKK25 Prop.~2.3; BKSW23 abstract (white-box 2nd pass); Beh23 \S1; ABR24 \S1 \\
Random ranks / RGMM (inside Beh21; phase-1 of Beh23/ABR24) & \textbf{yes} & ABR24 \S1: maximal matching ``using the fast algorithms of [BGS, Sol16, BDHSS19]'' \\
Vertex sparsification (additive$\to$multiplicative) & \textbf{yes} & BKSW23 Prop.~A.1 ($O(\log^2 n/\varepsilon^3)$ random contracted subgraphs, whp); AKK25 Prop.~2.4; BG24 Prop.~5; Beh23 \S1 \\
Adaptivity wrapper (fresh coins per rebuild; low-sensitivity outputs) & \textbf{yes} & BKSW23 abstract; BG24 \S3 + Alg.~2 ll.~17--20 + Lem.~9; AKK25 Result~1 + Problem~1 \\
\bottomrule
\end{tabular}
\end{center}

\noindent Two consequences, both load-bearing for WS-A:

\begin{enumerate}
\item \textbf{The deterministic frontier has exactly three modules.}
Everything above the midrule needs no derandomization. What remains is:
(a)~the opportunistic sampler (rows R4/R5 of Table~\ref{tab:audit});
(b)~the Beh21 estimation interface (R1/R2); (c)~vertex sparsification (R7).
\item \textbf{BG24's coins are narrower than the plan's v0.1 audit assumed.}
The ORS-epoch amortization itself (Lemma~11, $H_{\mathrm{cert}}$) is
deterministic combinatorics; the randomness of BG24 lives \emph{only} in
Lemma~7's sampling, plus the R1/R7 interfaces and the per-rebuild freshness.
This sharpens plan row R4 accordingly.
\end{enumerate}

\section{The audit table}\label{sec:table}

Table~\ref{tab:audit} is generated from the machine-readable source
(\texttt{auditlib/audit\_table.py}) by \texttt{run\_all.py}; the
\texttt{.tex} fragment below is \texttt{results/audit\_table.tex} verbatim.

\input{../results/audit_table.tex}

\section{The deterministic frontier, refined}\label{sec:frontier}

\paragraph{Module (b) is closed at the interface.} The estimation interface
consumes the kernel/contracted/leftover graph \emph{only through probes}
(adjacency-matrix queries in BG24 Prop.~4 / AKK25 Prop.~2.3). Plan
Lemma~3.6 shows that any \emph{deterministic} adaptive probe algorithm
distinguishing $\mu=0$ from $\mu\ge n/4$ needs more than
$\binom{\lceil n/2\rceil}{2}-1$ probes, against $\Otilde(n)$ randomized
(Prop.~2.3); plan Cor.~3.7 concludes there is no deterministic drop-in for
this layer, and experiment E5 (\S\ref{sec:experiments}) realizes the
adversary constructively. Derandomization must therefore \emph{replace} the
interface with maintained deterministic summaries (plan Tasks A1--A3), not
simulate it.

\paragraph{Module (a) is open --- and is now a precise target.} The
opportunistic sampler solves a \emph{search problem with a promise}:

\begin{problem}[Det-Opportunistic-Matching; refined Task A5 target]
\label{prob:detopp}
Given adjacency-matrix (BG24) or adjacency-list (AKK25) access to $G$, a set
$U\subseteq V$ with the promise $\mu(G[U])\ge\delta n$, and $\gamma<1/6$:
deterministically output a matching $M\subseteq G[U]$ with
$|M|\ge\gamma\delta n$ in time
$\Otilde\!\big(m\cdot n^{O(\gamma)}/\Delta_{\mathrm{IN}}(M)\big)$
(AKK25 shape; $\Otilde(n^{2+O(\gamma)}/d_{\mathrm{avg}})$ for the BG24
shape), where $\Delta_{\mathrm{IN}}(M)$ is the maximum degree of
$G[V(M)]$.
\end{problem}

\noindent Three observations delimit Problem~\ref{prob:detopp}:
(1)~it is \emph{not} covered by plan Lemma~3.6 --- that bound is for a
\emph{decision} task with no promise, whereas here a large matching is
promised to exist and the runtime is \emph{allowed} to degrade to
$\Otilde(m)$ exactly when the output is near-induced
($\Delta_{\mathrm{IN}}=O(1)$); (2)~the worst case $\Omega(m)$ is necessary
even for randomized algorithms (AKK25 quote ACK19), so only the
$d$-adaptivity needs replication; (3)~the randomized solutions get
$d$-adaptivity from one Chernoff-plus-union-bound argument (BG24 Claim~8;
AKK25 Claim~3.4 ``residual sparsity''), which suggests two deterministic
routes worth pricing in WS-A: degree-bucketed exhaustive scanning with a
potential-function charging in place of Claim~8, and expander/disperser-based
sample schedules (plan Task A4 \statusC). If Problem~\ref{prob:detopp} is
solved in the BG24 shape, the remaining obstacles to a deterministic
$\Otilde(\sqrt{n^{1+\varepsilon}\ORS})$ bound are exactly modules (b), (c)
--- this is the concrete content of plan Task A5's expectation that the
deterministic route lands at the BG24 shape first.

\paragraph{Related separations (prior-art check for the C1 note).}
Two published results already separate determinism from randomness for the
\emph{finding} task: Assadi--Solomon~\cite{AS19} (Thm.~3: deterministic
$\Omega(n^2)$ vs.\ randomized $\Otilde(n\beta)$ for finding a maximal
matching on adjacency-array graphs of neighborhood independence $\beta=2$)
and \cite{ACK19} (Thm.~6: finding a maximal matching needs $\Omega(n^2)$
queries even randomized). Plan Lemma~3.6 is a different statement ---
adjacency-\emph{matrix} probes, size-\emph{estimation} (where randomized is
$\Otilde(n)$~\cite{Beh21}, unlike finding), explicit clique-coverage
threshold, pipeline corollary --- but the C1 note must cite both prominently
and scope its novelty accordingly; the plan (v1.1, Task C1(a)) now records
this.

\paragraph{Limited independence: a verified obstruction (Task A4).}
Beh21's locality bound (Thm.~3.5: $\mathbb{E}[T(v,\pi)]=O(\bar d\log n)$
over uniform $v$ \emph{and} uniform $\pi$) is proved by the charging map
$\phi(\pi,\vec P)$ of Lemma~3.10, which rotates ranks along query paths ---
a bijection of the \emph{full} symmetric group. $k$-wise independent rank
families are not closed under these rotations, so the analysis does not
transfer to limited independence as-is. Experiment E3 nonetheless finds even
$2$-wise ranks empirically indistinguishable from full independence on
oblivious instances: Task A4 now has both a precise technical obstacle and an
experimental signal that the obstacle may be an artifact of the proof.
A second nuance (ABR24 \S2): in ABR24 the random order is additionally
load-bearing for the \emph{approximation} analysis of the $b$-matching
(even-spread-of-copies argument), not only for locality --- a derandomization
must replace that too.

\paragraph{A deterministic data point (ABR24 Thm.~2).} ABR24's two-pass
semi-streaming $(2-\sqrt2-\varepsilon)$-approximation is
\emph{deterministic}: the combinatorial core (maximal matching +
duplicate-free $b$-matching + blossom-inequality analysis) needs no coins
when the edges can be read in passes. The dynamic version's randomness is
forced purely by the probe access model --- exactly the localization that
Cor.~3.7 predicts.

\paragraph{Module (c).} Vertex sparsification (R7) is the same
plant-against-a-known-seed pattern one level up: a fixed contraction
partition can be defeated by an adversary routing the matching across
collision classes. Its operative statement is now pinned (BKSW23 Prop.~A.1:
$O(\log^2 n/\varepsilon^3)$ random contracted subgraphs per update, whp);
deterministic replacement candidates are degree-based (non-contractive)
vertex reductions.

\section{Experimental evidence (package \texttt{audit\_lab})}\label{sec:experiments}

All numbers: full scale, seed 20260611; reproduce with
\texttt{python3 run\_all.py --scale full}; tests T1--T6 cross-validate every
measured quantity (Blossom vs.\ bitmask DP; local oracle vs.\ global greedy;
construction validity). Mapping E$\to$rows is in Table~\ref{tab:audit}.

\begin{itemize}
\item \textbf{E1 (R1, what sampling buys).} $n{=}4000$, $\bar d{=}5$:
vertex-sampling estimate of $|\mathrm{RGMM}|$ has mean error
$0.0024\,n$ at $s{=}6400$ samples (scaling $\propto 1/\sqrt s$,
Fig.~\ref{fig:e1}), at $0.8$ oracle evaluations per vertex query;
Blossom-checked $|\mathrm{MM}|/\mu\in[0.805,1]$.
\item \textbf{E2 (R2, fixed ranks, locality).} On a rank-decreasing path
with $m{=}64{,}000$ edges, \emph{one} oracle query costs exactly $m$
evaluations; fully random ranks: mean $1.8$ (Fig.~\ref{fig:e2}).
\item \textbf{E3 (R2/A4, limited independence).} $k$-wise ranks with
$k\in\{2,3,4,8\}$ empirically match full independence in accuracy
(err$/n$: $0.0051$ vs.\ $0.0051$) and locality on oblivious instances.
Empirical only; consistent with the $\Theta(\log n)$-wise sufficiency
conjecture of rows R4/R5 \statusS.
\item \textbf{E4 (R2 dynamic).} Head-edge toggling on a $8000$-edge path:
fixed-rank greedy suffers $8000$ recourse per update (full flip);
random ranks, same oblivious updates: mean $1.20$ ($6667\times$).
Supports plan Task A3(b).
\item \textbf{E5 (R1/R4/R7, fixed samples; Lemma 3.6 companion).}
$n{=}96$: against a fixed pseudorandom probe set, the constructive
adversary plants $n/4$ matching edges the detector \emph{never} hits
(0 hits at every budget), while fresh random probes with budget
$B{=}1127=\binom{48}{2}{-}1$ detect with $p=0.993$. Extra finding: the
pseudorandom family still fails at $B$ equal to $90\%$ of all pairs ---
a fixed probe set is defeated by volume-free structure, and must instead
clique-cover $\lceil n/2\rceil$ vertices; randomness is being purchased
\emph{instead of} structured coverage.
\end{itemize}

\begin{figure}[h]
\centering
\begin{minipage}{0.48\textwidth}\centering
\includegraphics[width=\linewidth]{e1_concentration.png}
\caption{E1: estimator concentration.}\label{fig:e1}
\end{minipage}\hfill
\begin{minipage}{0.48\textwidth}\centering
\includegraphics[width=\linewidth]{e2_locality.png}
\caption{E2: locality, fixed vs.\ random ranks.}\label{fig:e2}
\end{minipage}
\end{figure}

\section{Status and next steps}\label{sec:next}

\textbf{Done (this pass).} Lemma-level locations for BG24 (complete for the
algorithmic part) and AKK25 (\S1--4.1); abstract/overview-level verification
for BKSW23, Beh23, ABR24; rows R1--R8 of Table~\ref{tab:audit}; experiments
E1--E5 with cross-validation; the three-module frontier and
Problem~\ref{prob:detopp}.

\textbf{Done in pass 2.} (1)~Beh21 internals read (\S1--3): the
rank-rotation obstruction for Task A4 is now precise; (2)~BKSW23
Prop.~A.1 pinned (vertex sparsification, R7); (3)~ABR24 adversary model
verified (Thm.~1: adaptive, worst-case) and \S1--4 read, yielding the
RGMM-for-approximation nuance and the deterministic-streaming data point;
(4)~\cite{ACK19} resolved (Thm.~6) and \cite{AS19} prior art folded into
the C1 scoping; (5)~the v0.3 table and these findings are folded into the
plan document, re-issued as v1.1 (Appendix~A there).

\textbf{Remaining (small).} BDHSS19/Sol16 maintenance internals (R2
residue); BKSW23 \S5+ robustness lemmas (R3 residue); Kis22 full text
(R7 residue --- operative form verified). These do not block WS-A start.

\begin{thebibliography}{99}\footnotesize
\bibitem{BKSW23} S.~Bhattacharya, P.~Kiss, T.~Saranurak, D.~Wajc.
Dynamic matching with better-than-2 approximation in polylogarithmic update
time. SODA 2023, pp.~100--128; J.~ACM 71(5):33, 2024. arXiv:2207.07438.
\bibitem{Beh23} S.~Behnezhad. Dynamic algorithms for maximum matching size.
SODA 2023, pp.~129--162. arXiv:2207.07607.
\bibitem{ABR24} A.~Azarmehr, S.~Behnezhad, M.~Roghani. Fully dynamic
matching: $(2-\sqrt2)$-approximation in polylog update time. SODA 2024,
pp.~3040--3061. arXiv:2307.08772.
\bibitem{BG24} S.~Behnezhad, A.~Ghafari. Fully dynamic matching and ordered
Ruzsa--Szemer\'edi graphs. FOCS 2024, pp.~314--327. arXiv:2404.06069 (v4).
\bibitem{AKK25} S.~Assadi, S.~Khanna, P.~Kiss. Improved bounds for fully
dynamic matching via ordered Ruzsa--Szemer\'edi graphs. SODA 2025,
pp.~2971--2990. arXiv:2406.13573.
\bibitem{Beh21} S.~Behnezhad. Time-optimal sublinear algorithms for matching
and vertex cover. FOCS 2021, pp.~873--884.
\bibitem{Kis22} P.~Kiss. Improving update times of dynamic matching
algorithms from amortized to worst case. ITCS 2022.
\bibitem{Tir18} S.~Tirodkar. Deterministic algorithms for maximum matching
on general graphs in the semi-streaming model. FSTTCS 2018.
\bibitem{McG05} A.~McGregor. Finding graph matchings in data streams.
APPROX--RANDOM 2005, pp.~170--181.
\bibitem{ACK19} S.~Assadi, Y.~Chen, S.~Khanna. Sublinear algorithms for
$(\Delta+1)$ vertex coloring. SODA 2019, pp.~767--786. (Thm.~6: finding a
maximal matching needs $\Omega(n^2)$ queries, even randomized; quoted in
AKK25 for the $\Omega(m)$ worst-case necessity of its Lemma~3.1.)
\bibitem{AS19} S.~Assadi, S.~Solomon. When algorithms for maximal
independent set and maximal matching run in sublinear time. ICALP 2019,
17:1--17:17. (Thm.~3: deterministic $\Omega(n^2)$ for finding a maximal
matching at neighborhood independence $\beta=2$.)
\bibitem{GP13} M.~Gupta, R.~Peng. Fully dynamic $(1+\varepsilon)$-approximate
matchings. FOCS 2013, pp.~548--557.
\bibitem{BLM20} S.~Behnezhad, J.~{\L}\k{a}cki, V.~Mirrokni. Fully dynamic
matching: beating 2-approximation in $\Delta^{\varepsilon}$ update time.
SODA 2020, pp.~2492--2508.
\bibitem{Plan} R.~Cao (prep.). Optimal approximation for fully dynamic
matching and ordered Ruzsa--Szemer\'edi graphs: a feasibility-checked
execution plan. Working document, June 2026 (Lemma~3.6, Cor.~3.7,
Tasks A0--A5).
\end{thebibliography}

\end{document}
